% Part: first-order-logic
% Chapter: introduction
% Section: sentences

\documentclass[../../../include/open-logic-section]{subfiles}

\begin{document}

\olfileid{fol}{int}{snt}

\olsection{\usetoken{P}{sentence}}

ښه، اوس د !!{structure}s او !!{formula}s دپاره د صدق (``پکښې
رښتونې'') د تعريف يوه طرحه لرو۔ خو دا يو اضافي څيز---د
!!a{variable} ګومارنه---غواړي، حال دا چې زموږ غوښتنه د
!!{sentence}s تعريف و۔ له ګومارنو څنګه ځان خلاصوو، او !!{sentence}s
څه دي؟

غالباً د صوري منطق له لومړۍ پېژندنې څخه مو د !!{variable}s او کميت
ټاکونکو د اړيکې بحث ياد وي۔ له کميت ټاکونکي څخه تل يو !!a{variable}
وروسته راځي، او بيا د !!{sentence} په هغه برخه کښې چې کميت ټاکونکے
ورباندې پلي کېږي (د هغه ``ساحه'')، !!{variable} د هماغه کميت ټاکونکي
له خوا ``تړلے'' ګڼو۔ په !!{formula}s کښې دا شرط نۀ و چې هر
!!{variable} دې ورته برابر کميت ټاکونکے ولري؛ بې له برابر کميت
ټاکونکي !!{variable}s ``ازاد'' يا ``ناتړلي'' دي۔ موږ به ټول هغه
!!{formula}s !!{sentence}s بولو چې ازاد !!{variable}s نۀ لري۔

په !!a{formula}~$!A$ کښې د !!a{variable} د يوې پېښې د تړل کېدو، د
هغې د تړونکي کميت ټاکونکي، او د هغې د ازاد پاتې کېدو وجداني مفکوره
بيا هم ګرانه نۀ ده۔ ښايي د دې د ازمېښت کومه طريقه مو زده کړې وي،
چې ښايي د قوسونو شمېرل هم پکښې وي۔ خو موږ بايد پر دقيق تعريف ټينګار
وکړو---او ځکه چې !!{formula}s مو په استقرا تعريف کړي، په
!!a{formula}~$!A$ کښې د !!a{variable}~$x$ ازادې او تړلې پېښې هم په
استقرا تعريفولے شو۔ مثلاً زموږ د ساده شوې ژبې دپاره داسې کېدے شي:
\begin{enumerate}
  \item که $!A$ اټومي وي، د~$x$ ټولې پېښې پکښې ازادې دي (يعنې په
  ~ $\Atom{\Obj P}{x}$ کښې د~$x$ پېښه ازاده ده)۔
  \item که $!A$ د~$\lnot !B$ بڼه ولري، نو په~$\lnot !B$ کښې د~$x$
  يوه پېښه هله او يوازې هله ازاده ده چې په~$!B$ کښې د~$x$ ورته
  پېښه ازاده وي (يعنې په~$!B$ کښې د متغيرونو ازادې پېښې کټ مټ په
  ~ $\lnot !B$ کښې ورته پېښې دي)۔
  \item که $!A$ د~$(!B \land !C)$ بڼه ولري، نو په~$(!B \land !C)$
  کښې د~$x$ يوه پېښه هله او يوازې هله ازاده ده چې په~$!B$ يا~$!C$
  کښې د~$x$ ورته پېښه ازاده وي۔
  \item که $!A$ د~$\lexists[x][!B]$ بڼه ولري، نو په~$!A$ کښې د~$x$
  هېڅ پېښه ازاده نۀ ده؛ او که د~$\lexists[y][!B]$ بڼه ولري، چې~$y$
  له~$x$ څخه بېل !!{variable} وي، نو په~$\lexists[y][!B]$ کښې د~$x$
  يوه پېښه هله او يوازې هله ازاده ده چې په~$!B$ کښې د~$x$ ورته
  پېښه ازاده وي۔
\end{enumerate}

کله چې د متغيرونو د ازادو او تړلو پېښو دقيق تعريف ولرو، په ساده
ډول ويلے شو: !!a{sentence} هر هغه !!{formula} ده چې د
!!{variable}s ازادې پېښې نۀ لري۔

\end{document}
